\documentclass[11pt,a4paper]{article}
\usepackage[utf8]{inputenc}
\usepackage[T1]{fontenc}
\usepackage[romanian]{babel}
\usepackage[margin=1.3cm]{geometry}
\usepackage{array}
\usepackage{tabularx}
\usepackage{enumitem}
\usepackage{graphicx}
\usepackage{setspace}
\usepackage{xcolor}
\usepackage{hyperref}
\usepackage{amssymb}
\setlength{\parindent}{0pt}
\newcommand{\fillline}[1][6cm]{\underline{\hspace{#1}}}
\newcommand{\fillbox}[1][0.4cm]{\framebox[#1]{\rule{0pt}{#1}}}
\newcommand{\checkboxempty}{$\square$}

\begin{document}

\begin{minipage}{0.3\textwidth}
\textbf{PRIMĂRIA}\\
\textbf{CLUJ-NAPOCA}
\end{minipage}%
\begin{minipage}{0.4\textwidth}
\centering
\textbf{*432008*}
\end{minipage}%
\begin{minipage}{0.3\textwidth}
\raggedleft
Nr. \fillline[3cm] / \fillline[1.5cm]
\end{minipage}

\vspace{0.3cm}
\hrule
\vspace{0.5cm}

\textbf{Către,}

\begin{center}
\textbf{\underline{PRIMĂRIA MUNICIPIULUI CLUJ-NAPOCA}}\\[0.3cm]
\textbf{CERERE}\\[0.2cm]
\textbf{PENTRU ELIBERAREA CERTIFICATULUI DE ATESTARE A}\\
\textbf{EDIFICĂRII/EXTINDERII CONSTRUCȚIEI}
\end{center}

\vspace{0.5cm}

\hspace{0.5cm}Subsemnatul \underline{ {{nume_complet}} } CNP/CUI \underline{ {{cnp}} }

cu domiciliul in județul \fillline[4cm], municipiul/orașul/comuna \fillline[5cm],

sectorul/satul \fillline[5cm], cod poștal \fillline[2cm], strada \fillline[6cm],

nr. \fillline[1.5cm], bl. \fillline[1cm], sc. \fillline[1cm], et. \fillline[1cm], ap. \fillline[1cm], telefon/fax \fillline[5cm]

\vspace{0.5cm}

\begin{center}
\textbf{solicit eliberarea}\\[0.2cm]
\textbf{în conformitate cu prevederile legale a CERTIFICATULUI DE ATESTARE A EDIFICĂRII/EXTINDERII CONSTRUCȚIEI, pentru construcțiile situate în județul CLUJ,}
\end{center}

\vspace{0.3cm}

municipiul CLUJ-NAPOCA, cartierul \fillline[7cm], cod poștal \fillline[2cm],

strada \underline{ {{adresa_imobil}} }, nr \fillline[3cm]

Nr. Fişa cadastrală/Nr. Carte Funciară \fillline[8cm]

\vspace{0.3cm}

Nr. topografic al parcelei \fillline[8cm]

\vspace{0.4cm}

\hspace{0.5cm}Anexez la prezenta cerere (în copie)* următoarele:

\vspace{0.3cm}

CERTIFICATUL DE URBANISM Nr. \fillline[8cm]

EXTRASUL DE CARTE FUNCIARĂ Nr. \fillline[8cm]

AUTORIZAȚIA DE CONSTRUIRE Nr. \underline{ {{nr_ac}} }

PROIECTUL ÎN FAZA P.A.C. AUTORIZAT Nr. \fillline[8cm]

\vspace{0.3cm}

Declar pe proprie răspundere că datele menționate în prezenta cerere sunt exacte.

\vspace{0.4cm}

Telefon: \fillline[5cm]

\vspace{0.4cm}

Data: \underline{ {{data_finalizare}} } \hfill Semnătura \fillline[5cm]

\vfill
\hrule
\vspace{0.2cm}
\small
www.primariaclujnapoca.ro \\
telefon: 0264/596.030 \\
Cluj-Napoca, str. Moţilor, nr. 3 \hfill \textbf{*432008*}

\vspace{0.2cm}
\footnotesize
Timp estimativ de completare: 10 minute\\
Datele personale care vă sunt solicitate prin prezenta cerere vor fi prelucrate numai în vederea procesării și soluționării solicitării dumneavoastră. Primăria municipiului Cluj-Napoca garantează securitatea procesării datelor și arhivarea acestora în conformitate cu prevederile legale în vigoare. Responsabilul Primăriei municipiului Cluj-Napoca cu protecția datelor poate fi contactat pe adresa de email dpo@primariaclujnapoca.ro.\\
În conformitate cu Regulamentul nr. 679/2016 aveți dreptul de a solicita Primăriei municipiului Cluj-Napoca, în ceea ce priveşte datele cu caracter personal referitoare la persoana vizată, accesul la acestea, rectificarea sau ştergerea acestora sau restricţionarea prelucrării sau a dreptului de a vă opune prelucrării, precum şi a dreptului la portabilitatea datelor.

\newpage

\begin{minipage}{0.3\textwidth}
\textbf{PRIMĂRIA}\\
\textbf{CLUJ-NAPOCA}
\end{minipage}%
\begin{minipage}{0.4\textwidth}
\centering
\textbf{*432008*}
\end{minipage}%
\begin{minipage}{0.3\textwidth}
\raggedleft
Nr. \fillline[3cm] / \fillline[1.5cm]
\end{minipage}

\vspace{0.3cm}
\hrule
\vspace{0.5cm}

\textbf{Actele necesare pentru emiterea certificatului de atestare a edificării/extinderii construcției:}

\begin{enumerate}[leftmargin=*]
\item Extrasul CF nu mai vechi de trei luni in original (cel electronic, numai vechi de 30 de zile);
\item Copie C.I. sau CUI;
\item Anuntul de incepere lucrari;
\item Fotografii datate/semnate si stampilate de proiectant si beneficiar;
\item Declaratia privind valoarea reala a lucrarilor (se va completa prima parte, fără Fișa de calcul);
\item Pentru persoane fizice -- Deviz general conform HG nr. 28/2008;
\item Pentru societati comerciale situatii de lucrari si adresa oficiala din partea firmei cu valoarea finala a lucrarilor cuprinsa in contabilitate;
\item Decizia de impunere de la taxe si impozite sau Proces verbal de stabilire a bazei impozabile pentru clădiri;
\item Proces-verbal de receptie la terminarea lucrarilor in 2 exemplare originale (anexa 2 la HG 343/2017);
\item Referatul proiectantului;
\item Referatul dirigintelui de santier;
\item Certificatul de performanta energetica a cladirii, conform Legii nr. 372/2005 privind performanta energetica a cladirilor, art. 12, alin. 1, ordinul nr. 1459/2007 si Ordinul ministrului dezvoltarii regionale si locuintei nr. 839/2009 pentru aprobara Normelor metodologice de aplicare a Legii nr. 50/1991 privind autorizarea executarii lucrarilor de constructii;
\item Adeverința eliberată de Inspectoratul de Stat in Constructii prin care se confirmă plățile efectuate, respectiv dacă investitorul a virat către I.S.C. sumele aferente cotelor prevăzute în Legea nr. 50/1991;
\item Dovada achitarii regularizarii taxei;
\item Dovada transportului materialului rezultat din lucrarile de constructii la rampe amenajate (factura sau chitanta);
\item Dovada bransarii constructiei la utilitati. Dovada poate fi facuta cu:
\begin{itemize}
\item pentru electrica: proces verbal de punere in functiune sau receptie a lucrarilor și aviz tehnic de racordare, contract de furnizare (definitiv) sau certificat de racordare;
\item pentru apă-canal: proces verbal de recepție tehnică sau contract de furnizare;
\item pentru gaze-naturale: contract de furnizare;
\end{itemize}
\item Plan de amplasament si delimitare a imobilului, intocmit de o persoana autorizata ANCPI (anexa 1.35 la Ordinul 700/2014 privind aprobarea Regulamentului de avizare, receptie si inscriere in evidentele de cadastru si carte funciara);
\item Autorizatia de construire insotita de documentatia tehnica (DT) vizata spre neschimbare in original.
\end{enumerate}

\vspace{0.4cm}

\scriptsize
Primăria municipiului Cluj-Napoca asigură gratuit fotocopierea documentelor.\\
Pentru documentele emise de către alte instituții publice, organe de specialitate ale administrației publice centrale și locale, precum și persoane juridice de drept privat, care potrivit legii au obținut statut de utilitate publică sau sunt autorizate să presteze un serviciu public, în regim de putere publică, pe care solicitantul nu le depune, acesta are posibilitatea de a-și da consimțământul expres ca Primăria municipiului Cluj-Napoca să obțină în numele său copii sau extrase ale acestora.

\vfill
\hrule
\vspace{0.2cm}
\small
www.primariaclujnapoca.ro \\
telefon: 0264/596.030 \\
Cluj-Napoca, str. Moţilor, nr. 3 \hfill \textbf{*432008*}

\end{document}
